\subsection{Local--global injectivity for the splitting transfer quotient}\label{sec:paper-Q-local-global}
In this subsection we prove the local--global principle for the splitting transfer quotient $Q_F(D)$ over global fields (Theorem~\ref{thm:paper-Q-local-global}) by comparing two global evaluations of a single universal $K_2$-class.
Proposition~\ref{prop:paper-global-picard} and Lemma~\ref{lem:paper-adelic-eval} are the global counterparts of the curve argument of Section~\ref{sec:paper-splitting}.
Here the norm equation must be solved simultaneously at all places $v\in S$.
Since only classes modulo $I_F(D)$ are needed, Lemma~\ref{lem:paper-curve-reciprocity}(1) handles the boundary terms, and no rigidification is required.

Throughout this subsection, $F$ is a global field and $D$ is a central division algebra over $F$.
Let $S$ be the finite set of places where $D$ does not split, including ramified real places, and write $D_v=D\otimes_FF_v$.

\begin{proposition}\label{prop:paper-global-picard}
Let $V$ be a nonempty smooth open subset of a projective geometrically integral curve $C/F$.
Let $P_D(C,V)\subset\operatorname{Pic}(C)$ be the subgroup generated by the divisor classes of the closed points of $V$ whose residue fields are separable over $F$ and split $D$.
Restriction induces an injection
\[
 \operatorname{Pic}(C)/P_D(C,V)\hookrightarrow
 \prod_{v\in S}\operatorname{Pic}(C_{F_v})/P_{D_v}(C_{F_v},V_{F_v}).
\]
\end{proposition}
\begin{proof}
Let $\mathcal L$ be a line bundle on $C$ whose class vanishes in every local quotient.
At each $v\in S$, we lift the separable split closed points to rational points over their residue fields and write
\[
 \mathcal L_{F_v}\simeq\bigotimes_iN_{E_{v,i}/F_v}(\mathcal M_{v,i})
\]
for line bundles $\mathcal M_{v,i}$ on $C_{E_{v,i}}$, where the $E_{v,i}/F_v$ are splitting fields of $D_v$; if this decomposition is empty, we add one trivial factor from a splitting field.
Put $n_v=\sum_i[E_{v,i}:F_v]$, choose a very ample line bundle $\mathcal H$ on $C$, and choose a sufficiently large common multiple $m$ of the $n_v$.
Twisting each $\mathcal M_{v,i}$ and each trivial bundle on $C_{E_{v,i}}$ by $\mathcal H^{\otimes m/n_v}$, we obtain
\[
 \bigotimes_iN_{E_{v,i}/F_v}\bigl(\mathcal M_{v,i}\otimes\mathcal H^{\otimes m/n_v}\bigr)\simeq\mathcal H^{\otimes m}\otimes\mathcal L,
 \qquad
 \bigotimes_iN_{E_{v,i}/F_v}\bigl(\mathcal H^{\otimes m/n_v}\bigr)\simeq\mathcal H^{\otimes m}.
\]

We next choose split divisors locally.
We take general sections of the twisted line bundles on the extensions $C_{E_{v,i}}$, and require their zero divisors to be geometrically reduced, to avoid the boundary $C\setminus V$, and to be disjoint from one another and from all their distinct conjugates, over all indices $i$ and all embeddings of the fields $E_{v,i}$.
Over $F_v^{\mathrm{sep}}$ the Weil restriction of the section space gives independent section spaces for the conjugates.
For two sections, vanishing at the same point imposes two independent linear conditions, so even when the point varies on the curve this is a proper closed condition.
Disjointness ensures that every point of a pushed-forward divisor has a residue field containing the corresponding splitting field, so we obtain split divisors in both localized complete linear systems $|\mathcal H^{\otimes m}\otimes\mathcal L|$ and $|\mathcal H^{\otimes m}|$ at $v$.

Finally we globalize.
Over the locus of divisors with simple zeros, the family of divisors is finite \'etale.
The inverse function theorem for the finite \'etale scheme of isomorphisms with a chosen local divisor then shows that nearby sections have the same local splitting type.
Weak approximation in the global section spaces gives two global divisors with the required conditions at every $v\in S$.
A closed point of either global divisor has a finite separable residue field $\kappa$ over $F$, and $D$ splits over every completion of $\kappa$; by the Brauer--Hasse--Noether theorem, $D_\kappa$ is split.
The difference of the two divisors represents $\mathcal L$.
When $S$ is empty, $D$ is split and general sections suffice.
\end{proof}

\begin{lemma}\label{lem:paper-adelic-eval}
Let $X/F$ be smooth quasi-projective and let $\alpha\in K_2(X;D)$.
For any $(b_v)\in\prod_{v\in S}X(F_v)$, the evaluation $[y^*\alpha]\in Q_F(D)$ is the same for all global points $y$ in a sufficiently small product neighborhood of $(b_v)$.
\end{lemma}
\begin{proof}
If the neighborhood contains a global point, we shrink it so that all its global points lie in a single connected component of $X$; this component is smooth and has an $F$-point, hence is geometrically integral, and we work there.
We fix a very ample line bundle on a projective closure of $X$; for a curve we use the fixed projective closure itself.
Dimension zero and the case $S=\varnothing$ are trivial.

We first construct curves through two nearby points.
In dimension at least two, we take \'etale coordinates
\[
 t_i=r_i/s_0
\]
at each $b_v$, where $r_i$ and $s_0$ are sections of the same sufficiently high power of the fixed very ample line bundle.
For every tangent direction $u$, we choose a geometrically integral complete-intersection curve through $b_v$ tangent to $u$; in sufficiently high degree, the sections satisfying the prescribed point and first-order conditions remain free away from the point, so Bertini's theorem applies.
We take the local curves to be cut out by sections of one sufficiently high degree, so that they form a projective flat family of geometrically integral curves over a common open set of parameters: the geometrically non-Cohen--Macaulay locus has codimension at least two, Bertini's theorem lets the curves avoid it, and the defining sections then form a regular sequence along them.
For the local norm we fix a finite separable splitting field.
The relative Picard norm is smooth, since it becomes addition after finite \'etale base change and $H^2(\mathcal O)=0$.
Since the curves have local rational points, Picard points are represented by line bundles.
By smoothness at the trivial solution on the diagonal, the condition that $\mathcal O(y-y')$ is a norm defines an open neighborhood of the diagonal in the family of curves with two smooth marked points $y,y'$.
We use this open neighborhood when the defining sections are moved by weak approximation.

Suppose that the secant direction of nearby points $y,y'$ is close to $u$.
Let $s_j$, $j\ge1$, be the defining sections of the central curve, and choose $i$ such that, for a fixed constant $c>0$,
\[
 |t_i(y')-t_i(y)|\ge c\,\|t(y')-t(y)\|.
\]
Sections vanishing at both points are given by
\[
 \beta_j=\frac{(s_j/s_0)(y')-(s_j/s_0)(y)}{t_i(y')-t_i(y)},\qquad
 s'_j=s_j-\bigl((s_j/s_0)(y)-\beta_j t_i(y)\bigr)s_0-\beta_jr_i.
\]
Since $d(s_j/s_0)(b_v)$ annihilates $u$, the correction $s'_j-s_j$ is small when the direction and the endpoints are close to $u$ and $b_v$, respectively; this follows by analytic expansion in \'etale coordinates over both real and nonarchimedean completions.
Compactness of the projective space of directions over $F_v$ gives a neighborhood of points that is uniform in the direction.

Now let $y,y'$ be global points in this neighborhood.
Weak approximation on the space of global sections vanishing at both points gives a single global curve, geometrically integral (an open condition preserved by approximation) and smooth at $y$ and $y'$, on which $\mathcal O(y-y')$ is a local norm from a splitting field at every $v\in S$.
As in the proof of Proposition~\ref{prop:paper-global-picard}, twisting by a very ample line bundle and pushing forward general sections shows that the local class of $\mathcal O(y-y')$ lies in the subgroup generated by split closed points.
Applied to the intersection $V$ of the smooth locus of the global curve with $X$, the proposition gives a divisor $Z$ supported on split closed points of $V$ with $[y-y']=[Z]$ in the Picard group.
Pulling back to the normalization gives $y-y'=\operatorname{div}(f)+Z$ for a rational function $f$.
Thus $\operatorname{div}(f)$ is supported in $V$, and $f$ is a unit at every branch of the boundary.
By Lemma~\ref{lem:paper-curve-reciprocity}(3),
\[
 y^*\alpha-y'^*\alpha\in I_F(D).
\]
\end{proof}


\begin{lemma}\label{lem:paper-local-I-single-field}
Let $k$ be a nonarchimedean local field and $B$ a central simple algebra over $k$, or let $k=\mathbb R$ and $B$ a nonsplit central simple algebra over $\mathbb R$.
For every finite separable splitting field $L/k$ of $B$,
\[
 I_k(B)=\theta_LK_2(L).
\]
\end{lemma}
\begin{proof}
Let $E/k$ be any finite splitting field of $B$, and put $M=EL$; the extension $M/E$ is separable.
For $k$ nonarchimedean, the field norm $N_{M/E}\colon K_2(M)\to K_2(E)$ is surjective \cite{BassTate}, \cite[p.~4]{BakRehmann1984}, so transitivity of transfer gives
\[
 \theta_E K_2(E)=\theta_M K_2(M)\subseteq\theta_LK_2(L).
\]
For $k=\mathbb R$, every splitting field is $\mathbb C$, and the same identity holds.
\end{proof}

For a unit $u$, we use the standard generators $x_{ij}(a)$ of the stable Steinberg group from Section~\ref{sec:paper-splitting} to define
\[
 w_{ij}(u)=x_{ij}(u)\,x_{ji}(-u^{-1})\,x_{ij}(u),\qquad h_{ij}(u)=w_{ij}(u)\,w_{ij}(-1).
\]
The matrix image of $h_{ij}(u)$ is diagonal with entries $u$ and $u^{-1}$ in positions $i$ and $j$, respectively, and $h_{ij}(1)=1$.

\begin{theorem}\label{thm:paper-Q-local-global}
For a global field $F$ and a central division algebra $D$ over $F$, the diagonal restriction
\[
 Q_F(D)\longrightarrow\prod_{v\in S}Q_{F_v}(D_v)
\]
is injective, where $S$ is the finite set of places where $D$ does not split.
\end{theorem}

For a class in the kernel, we lift its local images through the completions of one separable maximal subfield and compare two global evaluations of a $K_2$-class built from Steinberg words of a common shape.

\begin{proof}
Let $z\in K_2(D)$ have zero class in every local quotient.
If $S=\varnothing$, then $D$ is split, and there is nothing to prove.

We first show that the image of $z$ in the product of the local algebras comes from a separable maximal subfield.
Let $L\subset D$ be a separable maximal subfield, and put
\[
 D_S=\prod_{v\in S}D_v,\qquad L_S=L\otimes_F\prod_{v\in S}F_v=\prod_{v,\,w\mid v}L_w.
\]
For the base change $L_S\to D_S$ of $L\subset D$, the $(v,w)$-component of extension of scalars $K_2(L_S)\to K_2(D_S)$ is tensoring with the $(L_w,D_v)$-bimodule $e_wD_v$, where $e_w$ is the idempotent of the factor $L_w$.
The bimodule $e_wD_v$ is the base change to $L_w$ of the simple module $D$ over $L\otimes_F D^{\mathrm{op}}\simeq M_{\deg D}(L)$, so this component is $\theta_{L_w}$.
By Lemma~\ref{lem:paper-local-I-single-field}, $I_{F_v}(D_v)=\theta_{L_w}K_2(L_w)$ for every $v\in S$ and $w\mid v$.
Thus the image of $z$ in $K_2(D_S)$ comes from $K_2(L_S)$.

Let $\mathsf Z$ be a freely reduced Steinberg word representing $z$, and let $\mathsf T$ be a word in Matsumoto symbols representing a lift to $K_2(L_S)$ of the image of $z$ in $K_2(D_S)$.
We represent a symbol $\{a,b\}$ by $h_{12}(a)h_{12}(b)h_{12}(ab)^{-1}$ \cite{Matsumoto}, and a symbol in one factor of $L_S$ by taking the units equal to $1$ in the other factors.
Since $\mathsf Z=\mathsf T$ in the stable Steinberg group over $D_S$, we can write
\[
 \mathsf Z=\mathsf R\mathsf T
\]
in the free group on the elementary generators, where $\mathsf R$ is a finite product of conjugates of defining relations and their inverses.
We work at one matrix rank containing all the indices used, fix the indices and signs in $\mathsf R\mathsf T$, and use an independent variable for the coefficient of each occurrence of a letter.
Since $D_S$ is a product ring, this gives a single word shape at all $v\in S$, and the original coefficients give local parameter points $b_v$.

We construct two global points for this word shape.
In the first, we insert common approximating coefficients in the pairs of letters that cancel when $\mathsf R\mathsf T$ freely reduces to $\mathsf Z$, and we use the original coefficients of $\mathsf Z$ in the remaining letters; the value at this point is $z$.
In the second, we approximate the independent parameters of the defining relations, compute the sum and product coefficients after approximation, and approximate the units in $\mathsf T$ simultaneously at all $L_w$ by weak approximation on $L$; the value at this point lies in $\theta_L K_2(L)$.
At both points the matrix word is exactly the identity.

Finally, the additional letters of Lemma~\ref{lem:paper-steinberg-fiber} make the identity fiber smooth at all $b_v$, and the universal word defines a $K_2$-class on its smooth open subset.
If the two global points are close enough to the $b_v$, their evaluations $[z]$ and $0$ agree in $Q_F(D)$ by Lemma~\ref{lem:paper-adelic-eval}; hence $[z]=0$.
This proof does not use the vanishing of $\SK_2$ over local fields.
\end{proof}
